\documentclass[12pt]{article}
\usepackage[UTF8]{ctex}
\usepackage{amsmath, amssymb, amsthm}
\usepackage{geometry}
\usepackage{algorithm}
\usepackage{algpseudocode}
\geometry{margin=1in}

\newenvironment{solution}{\par\noindent\textbf{解：}\par}{\par}

\begin{document}

\begin{minipage}{\textwidth}
    \begin{center}
        \fbox{
            \begin{minipage}{0.95\textwidth}
                \begin{center}
                    {\large\textbf{ }}
                \end{center}
                \begin{center}
                    {\large\textbf{Problem Set 2}}
                \end{center}
                \begin{center}
                    {Due date: March 18, 2026}
                \end{center}
                \begin{center}
                    {倪伟丰 2024110089}
                \end{center}
                \begin{center}
                    {\large\textbf{ }}
                \end{center}
            \end{minipage}
        }
    \end{center}
\end{minipage}

\bigskip

\begin{enumerate}
    \item (Exercise 2.3 from the textbook) Take the following list of functions and arrange them in ascending order of growth rate. That is, if function $g(n)$ immediately follows function $f(n)$ in your list, then it should be the case that $f(n)$ is $O(g(n))$.

    \medskip
    \begin{center}
    $f_1(n) = n^{2.5}$

    $f_2(n) = \sqrt{2n}$

    $f_3(n) = n + 10$

    $f_4(n) = 10^n$

    $f_5(n) = 100^n$

    $f_6(n) = n^2 \log n$

    \end{center}

    \begin{solution}
        递增顺序为：
        $$
            f_2,\ f_3,\ f_6,\ f_1,\ f_4,\ f_5
        $$

        逐步说明如下：

        \begin{itemize}
            \item $f_2 = O(f_3)$：$f_2(n) = \sqrt{2n} = \Theta(\sqrt{n})$，而 $\sqrt{n} = O(n) = O(n + 10)$，故 $f_2 = O(f_3)$。

            \item $f_3 = O(f_6)$：$f_3(n) = n + 10 = \Theta(n)$，而 $n = O(n^2) = O(n^2 \log n)$，故 $f_3 = O(f_6)$。

            \item $f_6 = O(f_1)$：$\frac{f_6(n)}{f_1(n)} = \frac{n^2 \log n}{n^{2.5}} = \frac{\log n}{n^{0.5}} \to 0$（$n \to \infty$），故存在 $n_0$ 使得当 $n \geq n_0$ 时 $f_6(n) \leq f_1(n)$，即 $f_6 = O(f_1)$。

            \item $f_1 = O(f_4)$：对任意 $\varepsilon > 0$，当 $n$ 充分大时 $n^{2.5} \leq (1+\varepsilon)^n \leq 10^n$，故 $f_1 = O(f_4)$。

            \item $f_4 = O(f_5)$：$\frac{f_4(n)}{f_5(n)} = \frac{10^n}{100^n} = \left(\frac{1}{10}\right)^n \to 0$，故存在 $n_0$ 使得当 $n \geq n_0$ 时 $f_4(n) \leq f_5(n)$，即 $f_4 = O(f_5)$。
        \end{itemize}
    \end{solution}

    \bigskip

    \item (Exercise 2.5 from the textbook) Assume you have functions $f$ and $g$ such that $f(n)$ is $O(g(n))$. For each of the following statements, decide whether you think it is true or false and give a proof or counterexample.
    \begin{enumerate}
        \item $\log_2 f(n)$ is $O(\log_2 g(n))$.
        \item $2^{f(n)}$ is $O(2^{g(n)})$.
        \item $f(n)^2$ is $O(g(n)^2)$.
    \end{enumerate}

    \begin{enumerate}
        \item \textbf{错误。} 反例：取 $f(n) = 2$，$g(n) = 1$。

        取 $c = 2$，$n_0 = 1$，则对所有 $n \geq n_0$ 均有 $f(n) = 2 \leq 2 \cdot 1 = c \cdot g(n)$，故 $f(n) = O(g(n))$。

        但 $\log_2 f(n) = \log_2 2 = 1$，而 $\log_2 g(n) = \log_2 1 = 0$。

        由于 $\lim_{n \to \infty} \frac{\log_2 f(n)}{\log_2 g(n)} = \infty$，故 $\log_2 f(n)$ 不是 $O(\log_2 g(n))$。
  
        \item \textbf{错误。} 反例：取 $f(n) = 2n$，$g(n) = n$。

        取 $c = 2$，$n_0 = 1$，则对所有 $n \geq n_0$ 均有 $f(n) = 2n \leq 2n = c \cdot g(n)$，故 $f(n) = O(g(n))$。

        但 $\frac{2^{f(n)}}{2^{g(n)}} = \frac{2^{2n}}{2^n} = 2^n \to \infty$（$n \to \infty$）。

        因此不存在常数 $c' > 0$ 使得对所有充分大的 $n$ 均有 $2^{f(n)} \leq c' \cdot 2^{g(n)}$，故 $2^{f(n)}$ 不是 $O(2^{g(n)})$。

        \item \textbf{正确。} 因为 $f(n) = O(g(n))$，存在常数 $c > 0$ 和 $n_0 \geq 0$，使得对所有 $n \geq n_0$ 均有
        $$f(n) \leq c \cdot g(n)$$
        两边平方得
        $$f(n)^2 \leq c^2 \cdot g(n)^2$$
        取常数 $c' = c^2$，则对所有 $n \geq n_0$ 均有 $f(n)^2 \leq c' \cdot g(n)^2$，故 $f(n)^2 = O(g(n)^2)$。
    \end{enumerate}

    \bigskip

    \item (Exercise 2.6 from the textbook) Consider the following basic problem. You're given an array $A$ consisting of $n$ integers $A[1], A[2], \ldots, A[n]$. You'd like to output a two-dimensional $n$-by-$n$ array $B$ in which $B[i, j]$ (for $i < j$) contains the sum of array entries $A[i]$ through $A[j]$---that is, the sum $A[i] + A[i + 1] + \ldots + A[j]$. (The value of array entry $B[i, j]$ is left unspecified whenever $i \geq j$, so it doesn't matter what is output for these values.) Here's a simple algorithm to solve this problem.

    \medskip

    \begin{algorithmic}
    \For{$i = 1, 2, \cdots, n$}
        \For{$j = i + 1, i + 2, \cdots, n$}
            \State Add up array entries $A[i]$ through $A[j]$
            \State Store the result in $B[i, j]$
        \EndFor
    \EndFor
    \end{algorithmic}

    \medskip

    \begin{enumerate}
        \item For some function $f$ that you should choose, give a bound of the form $O(f(n))$ on the running time of this algorithm on an input of size $n$ (i.e., a bound on the number of operations performed by the algorithm).

        \item For this same function $f$, show that the running time of the algorithm on an input of size $n$ is also $\Omega(f(n))$. (This shows an asymptotically tight bound of $\Theta(f(n))$ on the running time.)

        \item Although the algorithm you analyzed in parts (a) and (b) is the most natural way to solve the problem---after all, it just iterates through the relevant entries of the array $B$, filling in a value for each---it contains some highly unnecessary sources of inefficiency. Give a different algorithm to solve this problem, with an asymptotically better running time. In other words, you should design an algorithm with running time $O(g(n))$, where $\lim_{n\to\infty} g(n)/f(n) = 0$.
    \end{enumerate}

    \begin{solution}

        \begin{enumerate}
            \item 取 $f(n) = n^3$，该算法的运行时间为 $O(n^3)$。

            对每一对满足 $1 \leq i < j \leq n$ 的下标，计算 $B[i,j] = A[i] + \cdots + A[j]$ 需要 $j - i$ 次加法，故总操作次数为
            $$
                T(n) = \sum_{i=1}^{n} \sum_{j=i+1}^{n} (j - i).
            $$
            令 $k = j - i$，则
            $$
                T(n) = \sum_{i=1}^{n-1} \sum_{k=1}^{n-i} k
                     = \sum_{i=1}^{n-1} \frac{(n-i)(n-i+1)}{2}
                     \leq \sum_{i=1}^{n-1} \frac{n^2}{2}
                     = \frac{(n-1)n^2}{2}
                     = O(n^3)
            $$

            \item 该算法的运行时间也是 $\Omega(n^3)$。

            考虑满足 $1 \leq i \leq \lfloor n/4 \rfloor$ 且 $\lceil 3n/4 \rceil \leq j \leq n$ 的下标对 $(i, j)$。这样的对至少有 $\lfloor n/4 \rfloor \cdot \lfloor n/4 \rfloor = \Omega(n^2)$ 个，且每个区间长度满足
            $$
                j - i \geq \left\lceil \frac{3n}{4} \right\rceil - \left\lfloor \frac{n}{4} \right\rfloor \geq \frac{n}{2}.
            $$
            因此
            $$
                T(n) \geq \Omega(n^2) \cdot \frac{n}{2} = \Omega(n^3).
            $$
            结合 (a) 的上界，得 $T(n) = \Theta(n^3)$。

            \item 利用前缀和可以将运行时间优化至 $O(n^2)$。算法如下：

        \end{enumerate}

        \begin{algorithmic}
            \State $S[0] \gets 0$
            \For{$k = 1, 2, \cdots, n$}
                \State $S[k] \gets S[k-1] + A[k]$
            \EndFor
            \For{$i = 1, 2, \cdots, n$}
                \For{$j = i+1, i+2, \cdots, n$}
                    \State $B[i, j] \gets S[j] - S[i-1]$
                \EndFor
            \EndFor
        \end{algorithmic}

        \textbf{正确性：} 由前缀和定义 $S[k] = A[1] + \cdots + A[k]$，故
        $$
            S[j] - S[i-1] = A[i] + A[i+1] + \cdots + A[j] = B[i,j].
        $$

        \textbf{时间复杂度：} 构造前缀和数组 $S$ 需要 $O(n)$ 次操作；双重循环共执行 $\binom{n}{2} = \Theta(n^2)$ 次，每次仅需一次减法（$O(1)$）。故总运行时间为 $O(n^2)$。

        由于 $\lim_{n \to \infty} \dfrac{n^2}{n^3} = 0$，该算法渐近复杂度严格优于原算法。

    \end{solution}
\end{enumerate}

\end{document}
